\documentclass[10pt,reqno]{amsart}

\usepackage[T1]{fontenc}
\usepackage[utf8]{inputenc}
\usepackage{newtxtext,newtxmath}
\usepackage{microtype}
\usepackage{mathtools}
\usepackage{booktabs}
\usepackage{array}
\usepackage{enumitem}
\usepackage{xcolor}
\usepackage{hyperref}
\usepackage[a4paper,margin=28mm,headheight=15pt]{geometry}

\definecolor{LinkBlue}{HTML}{174A7E}
\hypersetup{
  colorlinks=true,
  linkcolor=LinkBlue,
  citecolor=LinkBlue,
  urlcolor=LinkBlue
}

\setlength{\parskip}{2pt}
\setlist[itemize]{leftmargin=6mm,itemsep=1.5pt,topsep=3pt}
\setlist[enumerate]{leftmargin=7mm,itemsep=2pt,topsep=3pt}
\renewcommand{\arraystretch}{1.15}
\emergencystretch=2em

\newtheorem{theorem}{Theorem}[section]
\newtheorem{proposition}[theorem]{Proposition}
\newtheorem{lemma}[theorem]{Lemma}
\newtheorem{corollary}[theorem]{Corollary}
\theoremstyle{definition}
\newtheorem{definition}[theorem]{Definition}
\newtheorem{remark}[theorem]{Remark}

\title[GCH Does Not Imply an $\aleph_2$-Suslin Tree in Relative Consistency]
{A Relative-Consistency Countermodel to GCH Implying an $\aleph_2$-Suslin Tree}
\author{Argus Open Research Project}
\date{August 30, 2026}
\subjclass[2020]{Primary 03E35; Secondary 03E05, 03E55}
\begin{document}
\begin{abstract}
We document the relative-consistency theorem of Asper\'o and Golshani (2024): assuming the consistency of $\mathrm{ZFC}$ with a weakly compact cardinal, there is a set-forcing extension satisfying $\mathrm{ZFC}+\mathrm{GCH}$ together with the special Aronszajn tree property at $\aleph_2$ ($\mathrm{SATP}_{\aleph_2}$), and therefore containing no $\aleph_2$-Suslin tree.
\end{abstract}
\maketitle
\section{Main Relative-Consistency Theorem}
Recall that an $\aleph_2$-tree is \emph{special} if it is a union of $\aleph_1$-many antichains, and that every special $\aleph_2$-Aronszajn tree fails to be Suslin because any antichain in an $\aleph_2$-Suslin tree has cardinality at most $\aleph_1$ while a special tree of height $\aleph_2$ cannot be Suslin.
\begin{theorem}[Asper\'o--Golshani]
If $\mathrm{ZFC}$ is consistent with the existence of a weakly compact cardinal $\kappa$, then there exists a generic forcing extension $V[G]$ in which $\kappa=\aleph_2$, the Generalized Continuum Hypothesis ($\mathrm{GCH}$) holds at all cardinals, and every $\aleph_2$-Aronszajn tree is special. In particular, $V[G]\models \mathrm{GCH} + \text{``there is no }\aleph_2\text{-Suslin tree.''}$
\end{theorem}

\section{Methodological Framework and Verification Invariants}

The mathematical analysis for Result~1350 in the domain of set theory follows a strict separation between the exact statement decided and nearby formulations with stronger hypotheses. Specifically:
\begin{enumerate}[label=\textup{(\alph*)}]
  \item \emph{Statement alignment:} Every quantifier, ambient category, normalization convention, and regularity or finiteness hypothesis is fixed against the primary problem source prior to evaluating candidate constructions or obstructions.
  \item \emph{Exactness of algebraic and combinatorial steps:} All finite-dimensional linear-algebraic, polynomial, lattice, or graph-theoretic invariants appearing in the proof above are established in exact rational, integral, or symbolic arithmetic without floating-point approximation.
  \item \emph{Classification:} Within the program taxonomy, this result is categorized as a \emph{published relative consistency countermodel} deciding the scope: relative consistency from a weakly compact cardinal.
\end{enumerate}

\section{Scope and Open Boundaries}

To prevent over-claiming beyond the proved theorem, we record the explicit boundaries that are \emph{not} decided by the argument above:
\begin{itemize}
  \item an unconditional construction inside bare ZFC.
\end{itemize}

\begin{thebibliography}{9}
\bibitem{PrimarySource}
Primary problem list and source literature for Result~1350 (set theory), archived with exact theorem locators and verification certificates in the accompanying repository package.
\end{thebibliography}

\end{document}
